% BHU PYQ -- Manually transcribed & error-corrected
\documentclass[11pt,a4paper]{article}

\usepackage[a4paper,top=2.5cm,bottom=2.5cm,left=2.8cm,right=2.8cm,headheight=14pt]{geometry}
\usepackage{amsmath,amssymb}
\usepackage[T1]{fontenc}
\usepackage[utf8]{inputenc}
\usepackage{lmodern}
\usepackage{microtype}
\usepackage{fancyhdr}
\usepackage{enumitem}
\usepackage{hyperref}
\usepackage{lastpage}
\usepackage{parskip}

\hypersetup{colorlinks=true,urlcolor=black,linkcolor=black,
  pdftitle={CHB-301: Organic & Physical Chemistry-II -- BHU PYQ},pdfauthor={Banaras Hindu University}}

\pagestyle{fancy}\fancyhf{}
\renewcommand{\headrulewidth}{0.4pt}\renewcommand{\footrulewidth}{0.4pt}
\fancyhead[L]{\small\textit{Banaras Hindu University}}
\fancyhead[R]{\small\textit{CHB-301: Organic \& Physical Chemistry-II}}
\fancyfoot[C]{\small Page \thepage\ of \pageref{LastPage}}

\newlist{parts}{enumerate}{2}
\setlist[parts,1]{label=(\alph*),leftmargin=*,itemsep=5pt,topsep=3pt}
\setlist[parts,2]{label=(\roman*),leftmargin=*,itemsep=3pt,topsep=2pt}
\newcommand{\pts}[1]{\hfill\textit{[#1]}}

\begin{document}

\begin{center}
  {\large\bfseries BANARAS HINDU UNIVERSITY}\\[2pt]
  {\normalsize B.Sc. (Hons.) Semester III Examination, 2023-24}\\[6pt]
  \rule{0.8\linewidth}{0.5pt}\\[6pt]
  {\large\bfseries Paper No. CHB-301: Organic \& Physical Chemistry-II}\\[4pt]
  {\normalsize Time: 3 Hours\hfill Maximum Marks: 70}
\end{center}
\rule{\linewidth}{0.4pt}
\small
\noindent\textit{Instructions: Answer five questions in total, including Question~1 from each section which is compulsory. Use separate answer books for Section~A and Section~B.}
\normalsize
\bigskip

\begin{center}
  {\large\bfseries SECTION A}\\[2pt]
  {\normalsize (Organic Chemistry-II -- Marks: 35)}
\end{center}
\rule{0.4\linewidth}{0.2pt}
\smallskip

\noindent\textbf{Question 1.} Answer all parts of the following: \pts{5 $\times$ 2 = 10}
\begin{parts}
  \item Classify the following species as aromatic, antiaromatic, or nonaromatic: (i) Cyclopropenyl anion, (ii) Cyclopentadienyl cation, (iii) Tropylium ion, and (iv) $1,3,5,7$-cyclooctatetraene.
  \item Identify structures A and B along with the mechanism of their formation:
    \[ \text{Ethylene oxide} \xrightarrow{\text{CH}_3\text{MgBr, THF}} A \xrightarrow{\text{H}_2\text{O}} B \]
  \item Arrange the following compounds in decreasing order of acid strength, with suitable reasons: phenol, $p$-nitrophenol, $p$-cresol, and $m$-nitrophenol.
  \item Draw the molecular orbital picture of benzene, indicating all nodal planes.
  \item What is DDT? Write down the reaction mechanism of its preparation from chloral and chlorobenzene.
\end{parts}

\medskip\rule{\linewidth}{0.2pt}

\noindent\textbf{Question 2.}
\begin{parts}
  \item Discuss the following reactions along with mechanisms: \pts{3 + 3 = 6}
  \begin{parts}
    \item Beckmann rearrangement
    \item Baeyer-Villiger oxidation
  \end{parts}
  \item Provide a suitable mechanism for the following reaction: \pts{3}
    \[ \text{Benzaldehyde} + \text{Acetic anhydride} \xrightarrow{\text{CH}_3\text{COONa}} \text{Cinnamic acid} \]
  \item Identify the products (A--G) in the following reaction sequence: \pts{3}
    \[ \text{Benzene} \xrightarrow{\text{CO, HCl, AlCl}_3, \text{CuCl}} A \xrightarrow{\text{dil. NaOH}} B \xrightarrow{\text{H}_3\text{O}^+} C \]
\end{parts}

\medskip\rule{\linewidth}{0.2pt}

\noindent\textbf{Question 3.}
\begin{parts}
  \item Explain why acetanilide undergoes mainly $p$-electrophilic substitution reactions, while aniline in strongly acidic medium exclusively produces $m$-products. \pts{6}
  \item Describe the synthesis of Novolac and Bakelite from phenol, including reactions and mechanisms. \pts{6}
\end{parts}

\medskip\rule{\linewidth}{0.2pt}

\noindent\textbf{Question 4.}
\begin{parts}
  \item An organic compound (A) with molecular formula $\text{C}_7\text{H}_8\text{O}$ reacts with neutral $\text{FeCl}_3$ to give a violet color. A reacts with $\text{CO}_2$ and $\text{NaOH}$ at high pressure to give B, which on acidification yields C. Identify A, B, and C, and write the mechanisms involved. \pts{6}
  \item Identify the products from the following reactions and discuss the mechanism of their formation: \dots \pts{6}
  \begin{parts}
    \item $\text{Pyridine} \xrightarrow{\text{NaNH}_2, 100^\circ\text{C}} ?$
    \item $\text{Pyrrole} \xrightarrow{\text{CHCl}_3, \text{KOH}} ?$
  \end{parts}
\end{parts}

\newpage
\begin{center}
  {\large\bfseries SECTION B}\\[2pt]
  {\normalsize (Physical Chemistry-II -- Marks: 35)}
\end{center}
\rule{0.4\linewidth}{0.2pt}
\smallskip

\noindent\textbf{Question 5.} Answer all parts of the following: \pts{5 $\times$ 2 = 10}
\begin{parts}
  \item Construct a cell from which one can measure the standard reduction potential of the $\text{Ag}^+(\text{aq}) + \text{e}^- \to \text{Ag(s)}$ half-cell reaction.
  \item Transference numbers of ions depend on concentration and temperature. Explain.
  \item The equivalent conductivity of $\text{CH}_3\text{COONa}$, $\text{HCl}$, and $\text{NaCl}$ at infinite dilution at $25^\circ\text{C}$ are given as $91.0$, $426.2$, and $126.5\,\text{ohm}^{-1}\,\text{cm}^2\,\text{equiv}^{-1}$, respectively. Calculate the equivalent conductivity of acetic acid at infinite dilution.
  \item Draw with brief explanation the approximate curve for the conductometric titration of $\text{NH}_4\text{OH}$ solution with $\text{HBr}$ solution as titrant.
  \item Under what conditions does the Gibbs phase rule become $F = C - P + 1$?
\end{parts}

\medskip\rule{\linewidth}{0.2pt}

\noindent\textbf{Question 6.}
\begin{parts}
  \item Show that for a solution of a single electrolyte, the mean ionic activity ($a_\pm$) is related to its molality ($m$) by $a_\pm = \gamma_\pm m_\pm$. Calculate $a_\pm$ for a $0.01\,\text{mol/kg}$ solution of $\text{Al}_2(\text{SO}_4)_3$. \pts{6}
  \item Describe the construction and working of the Quinhydrone electrode. What are its advantages and limitations? \pts{6}
\end{parts}

\medskip\rule{\linewidth}{0.2pt}

\noindent\textbf{Question 7.}
\begin{parts}
  \item The e.m.f. of a cell consisting of a hydrogen electrode and SCE was found to be $0.682\,\text{V}$ at $25^\circ\text{C}$ in a buffer solution. The e.m.f. increased to $0.794\,\text{V}$ when another buffer solution was used. What was the pH of the second buffer? \pts{6}
  \item How do buffer solutions resist changes in pH upon addition of small amounts of acid or base? Derive the Henderson-Hasselbalch equation for an acidic buffer. \pts{6}
\end{parts}

\medskip\rule{\linewidth}{0.2pt}

\noindent\textbf{Question 8.}
\begin{parts}
  \item Describe the procedure to determine the equilibrium constant of the reaction $\text{KI} + \text{I}_2 \rightleftharpoons \text{KI}_3$ in aqueous medium using the Nernst distribution law. \pts{6}
  \item Derive Bragg's equation for X-ray diffraction. If the first-order diffraction from a crystal plane occurs at a glance angle of $30^\circ$ using $\text{Cu-K}_\alpha$ radiation ($\lambda = 1.54\,\text{\AA}$), calculate the interplanar spacing ($d$). \pts{6}
\end{parts}

\vfill
\rule{\linewidth}{0.4pt}
\begin{center}\small
  \href{https://bhu-exam.vercel.app/}{Click here to visit our website}\\[4pt]
  \href{https://me-aryan.vercel.app/}{Click here to visit our contributor}\\[4pt]
  \href{https://quashan-me.vercel.app/}{Click here to visit our contributor}
\end{center}

\end{document}
